\section{Gabarito e resoluções comentadas — Legislação Específica do TCE-MA}

\begin{solucao}{07.01}\textbf{Errado.} O auxílio ao Poder Legislativo não cria relação hierárquica. O Tribunal exerce competências constitucionais próprias e possui autonomia funcional para deliberar nos limites de sua jurisdição.\end{solucao}
\begin{solucao}{07.02}\textbf{Certo.} Regimento disciplina organização e procedimento, mas não pode contrariar lei. A leitura deve respeitar Constituição, Lei Orgânica e atos infralegais em sua hierarquia.\end{solucao}
\begin{solucao}{07.03}\textbf{Gabarito: B.} Contas de governo são apreciadas mediante parecer prévio; contas de responsáveis submetidos à jurisdição de contas podem ser julgadas pelo Tribunal. As demais alternativas absolutizam ou eliminam distinções essenciais.\end{solucao}
\begin{solucao}{07.04}\textbf{Certo.} Instrução técnica subsidia a decisão. Relatório técnico, parecer ministerial, voto/proposta e deliberação são atos com funções distintas.\end{solucao}
\begin{solucao}{07.05}\textbf{Gabarito: B.} Débito está ligado ao ressarcimento do dano; multa é sanção. Podem coexistir, mas não se confundem.\end{solucao}
\begin{solucao}{07.06}\textbf{Certo.} A lógica da tomada de contas especial é reconstruir fato, autoria/responsabilidade e dano, permitindo decisão e eventual ressarcimento.\end{solucao}
\begin{solucao}{07.07}\textbf{Certo.} A IN nº 82/2025 foi editada para fiscalização da execução de emendas parlamentares estaduais e municipais, enfatizando transparência, rastreabilidade e conformidade.\end{solucao}
\begin{solucao}{07.08}\textbf{Gabarito: B.} A Lei nº 9.936/2013 aparece no edital como diploma de organização administrativa do TCE-MA.\end{solucao}
\begin{solucao}{07.09}\textbf{Certo.} Cautelar protege o resultado útil do controle ou evita dano; sanção é consequência de responsabilização e exige base legal e devido processo.\end{solucao}
\begin{solucao}{07.10}\textbf{Certo.} O Ministério Público de Contas possui função institucional própria e não é unidade técnica de auditoria.\end{solucao}

\begin{solucao}{07.11}\textbf{Errado.} A manifestação técnica não substitui a decisão da autoridade ou colegiado competente.\end{solucao}
\begin{solucao}{07.12}\textbf{Gabarito: B.} Decisão definitiva gravosa exige contraditório e defesa; situação urgente pode admitir cautelar com contraditório diferido quando juridicamente cabível.\end{solucao}
\begin{solucao}{07.13}\textbf{Certo.} O Regimento distribui competências entre relator e órgãos colegiados conforme matéria e fase processual.\end{solucao}
\begin{solucao}{07.14}\textbf{Gabarito: A.} Recurso deve ser memorizado como um conjunto: cabimento, legitimidade, prazo, efeito e competência. Decorar somente prazo é fonte clássica de erro.\end{solucao}
\begin{solucao}{07.15}\textbf{Certo.} A classificação das contas depende da natureza, gravidade, materialidade e efeitos da irregularidade; não decorre mecanicamente de qualquer falha formal.\end{solucao}
\begin{solucao}{07.16}\textbf{Certo.} Quantificar dano sem demonstrar responsável e nexo é insuficiente para imputação patrimonial.\end{solucao}
\begin{solucao}{07.17}\textbf{Gabarito: B.} A tomada de contas especial possui finalidade de apuração e ressarcimento, e não funções legislativas, penais ou orçamentárias.\end{solucao}
\begin{solucao}{07.18}\textbf{Certo.} A IN nº 50/2017 atribui papel à autoridade administrativa competente na instauração do procedimento, conforme o caso.\end{solucao}
\begin{solucao}{07.19}\textbf{Certo.} A IN nº 82/2025 foi expressamente vinculada pelo próprio TCE-MA ao ambiente jurídico criado pelas decisões na ADPF 854 e à necessidade de transparência das emendas.\end{solucao}
\begin{solucao}{07.20}\textbf{Gabarito: B.} Sem identificação do proponente, objeto, plano, cronograma e execução não é possível reconstruir a cadeia do recurso.\end{solucao}
\begin{solucao}{07.21}\textbf{Certo.} Transparência útil ao controle precisa acompanhar planejamento, execução orçamentária/financeira e entrega, não apenas o valor autorizado.\end{solucao}
\begin{solucao}{07.22}\textbf{Gabarito: B.} A lei prevalece sobre regra regimental incompatível. O intérprete deve harmonizar o procedimento ao diploma superior.\end{solucao}
\begin{solucao}{07.23}\textbf{Certo.} Alterações regimentais devem observar o procedimento formal próprio.\end{solucao}
\begin{solucao}{07.24}\textbf{Errado.} Parecer prévio e julgamento de contas são deliberações de natureza e destinatários institucionais distintos.\end{solucao}
\begin{solucao}{07.25}\textbf{Gabarito: B.} Rastreabilidade exige encadear origem, autorização, execução financeira, documentação e resultado físico.\end{solucao}

\begin{solucao}{07.26}\textbf{Errado.} Autonomia regimental não autoriza criar sanção sem base legal. Regimento não substitui reserva legal aplicável à sanção.\end{solucao}
\begin{solucao}{07.27}\textbf{Gabarito: B.} Débito requer dano, responsável, conduta e nexo sustentados por evidência. O valor elevado não elimina esse ônus.\end{solucao}
\begin{solucao}{07.28}\textbf{Certo.} Medida cautelar pode anteceder a oitiva quando a urgência justificar, sem equivaler a decisão definitiva de mérito.\end{solucao}
\begin{solucao}{07.29}\textbf{Certo.} Legislação específica é vulnerável a alterações. O mesmo nome de instituto pode permanecer enquanto prazo, competência ou efeito mudam.\end{solucao}
\begin{solucao}{07.30}\textbf{Gabarito: A.} O Regimento é a fonte central para distribuição interna de competência processual, lido em conjunto com a Lei Orgânica.\end{solucao}
\begin{solucao}{07.31}\textbf{Errado.} Total anual sem cadeia de execução não satisfaz rastreabilidade.\end{solucao}
\begin{solucao}{07.32}\textbf{Gabarito: B.} A fase executória precisa ser transparente: medição, documento fiscal, pagamento e entrega são elementos necessários à verificação.\end{solucao}
\begin{solucao}{07.33}\textbf{Certo.} O TCE-MA descreveu a IN nº 82/2025 como norma que afeta o ciclo completo da política pública financiada por emendas.\end{solucao}
\begin{solucao}{07.34}\textbf{Certo.} Lei de organização administrativa e Regimento possuem objetos distintos, ainda que se complementem.\end{solucao}
\begin{solucao}{07.35}\textbf{Gabarito: B.} Diligência integra instrução e busca formar evidência antes da decisão de mérito.\end{solucao}
\begin{solucao}{07.36}\textbf{Certo.} Defesa efetiva precisa poder questionar fato, metodologia, dano, nexo e autoria.\end{solucao}
\begin{solucao}{07.37}\textbf{Errado.} Recomendação anterior não elimina análise de fatos novos nem dispensa evidência em exercício posterior.\end{solucao}
\begin{solucao}{07.38}\textbf{Gabarito: A.} A sequência representa a lógica probatória da tomada de contas especial.\end{solucao}
\begin{solucao}{07.39}\textbf{Certo.} A IN nº 50/2017 trata de tomada de contas especial; a IN nº 82/2025, de emendas parlamentares.\end{solucao}
\begin{solucao}{07.40}\textbf{Gabarito: A.} Fiscalização é gênero e pode envolver múltiplos instrumentos e produtos, não apenas julgamento de contas.\end{solucao}

\begin{solucao}{07.41}\textbf{Errado.} Campos completos não bastam se a integridade histórica puder ser manipulada. Rastreabilidade exige trilha confiável de alterações.\end{solucao}
\begin{solucao}{07.42}\textbf{Gabarito: A.} Sem relação entre pagamentos, fornecedores, documentos e metas não se demonstra como o recurso produziu o resultado declarado.\end{solucao}
\begin{solucao}{07.43}\textbf{Certo.} Publicidade é condição necessária, mas rastreabilidade exige reconstruir a cadeia do recurso.\end{solucao}
\begin{solucao}{07.44}\textbf{Gabarito: B.} TCE especial não deve ser banalizada como procedimento para qualquer falha; seu objeto é apuração de fatos, responsáveis e dano nas hipóteses próprias.\end{solucao}
\begin{solucao}{07.45}\textbf{Certo.} Cargo ou posição hierárquica não substituem prova de conduta e nexo de responsabilidade.\end{solucao}
\begin{solucao}{07.46}\textbf{Gabarito: B.} Controle de privilégios, versionamento e trilha de auditoria preservam integridade e responsabilização sobre mudanças em dados sensíveis.\end{solucao}
\begin{solucao}{07.47}\textbf{Certo.} Fiscalização de emenda exige cruzar dimensões jurídicas, orçamentárias, financeiras e físicas.\end{solucao}
\begin{solucao}{07.48}\textbf{Gabarito: B.} Evidência de baixa confiabilidade não sustenta imputação robusta. É necessário reconstruir fonte, integridade e método, além de assegurar defesa.\end{solucao}
\begin{solucao}{07.49}\textbf{Certo.} A disciplina do caso pode estar distribuída em diplomas distintos, cuja leitura deve respeitar hierarquia e especialidade.\end{solucao}
\begin{solucao}{07.50}\textbf{Gabarito: B.} Os elementos indicam fragilidade simultânea de planejamento, documentação e rastreabilidade, impedindo comprovar adequadamente a regular aplicação do recurso.\end{solucao}
